\documentclass[12pt, a4paper]{article}

% ---- Standard Packages ----
\usepackage[a4paper, margin=1in]{geometry}
\usepackage[utf8]{inputenc}
\usepackage[T1]{fontenc}
\usepackage[bahasa]{babel}
\usepackage{amsmath, amssymb, amsthm, mathtools}
\usepackage{enumitem}
\usepackage{booktabs}
\usepackage{multicol}
\usepackage{hyperref}

\title{MODUL BELAJAR IPAS KELAS 10\\ \large Persiapan Ulangan Akhir Semester}
\author{Ringkasan Materi Komprehensif}
\date{}

\begin{document}
\maketitle

\tableofcontents
\newpage

\section{Aspek Keruangan, Geografi, dan Konektivitas Antarruang}

\subsection{Hakikat Ruang dan Pemahaman Dasar Alat Bantu Geografi}
Dalam ilmu geografi, ruang adalah tempat di permukaan bumi, baik secara keseluruhan maupun hanya sebagian, yang digunakan oleh makhluk hidup untuk tinggal. Setiap ruang di permukaan bumi memiliki karakteristik fisik dan sosial yang berbeda-beda. Untuk memahami karakteristik tersebut, manusia menciptakan berbagai alat bantu pemetaan.

\subsubsection{Pengertian dan Jenis-Jenis Peta}
Peta merupakan gambaran sebagian atau seluruh permukaan bumi pada suatu bidang datar yang diperkecil dengan menggunakan skala tertentu, serta dilengkapi dengan tulisan dan simbol sebagai keterangan. Peta berfungsi untuk menunjukkan lokasi, arah, jarak, luas, dan bentuk permukaan bumi.

Secara umum, peta dibagi menjadi dua jenis utama, yaitu:
\begin{enumerate}
    \item \textbf{Peta Umum:} Peta yang menggambarkan kenampakan bumi secara keseluruhan, baik kenampakan alam (pegunungan, sungai) maupun kenampakan buatan manusia (jalan, kota). Contoh peta umum adalah peta topografi (menggambarkan relief/tinggi rendah permukaan bumi) dan peta korografi.
    \item \textbf{Peta Tematik/Khusus:} Peta yang hanya menggambarkan informasi dengan tema tertentu. Contohnya: peta kepadatan penduduk, peta curah hujan, peta geologi, dan peta jaringan kabel serat optik bawah laut (sangat relevan dengan bidang TKJ).
\end{enumerate}

\subsubsection{Pengertian Atlas}
Atlas adalah sekumpulan peta yang dirangkai, disusun, dan dibukukan menjadi satu kesatuan. Atlas biasanya tidak hanya berisi peta, tetapi juga dilengkapi dengan data statistik, gambar, profil negara atau wilayah, daftar indeks letak suatu kota, serta penjelasan letak astronomis dan geografis. Berdasarkan jangkauannya, atlas dapat berupa atlas nasional, atlas regional (misalnya Atlas Asia Tenggara), maupun atlas dunia.

\subsubsection{Jenis Tiruan Bola Bumi (Globe)}
Globe adalah model tiruan bola bumi dalam bentuk tiga dimensi yang dibuat sedemikian rupa sehingga paling mendekati bentuk bumi yang sesungguhnya. Dibandingkan dengan peta, globe memiliki keunggulan mutlak: globe meminimalisir distorsi (penyimpangan) bentuk, luas, jarak, dan arah kawasan. 
Globe sering digunakan untuk:
\begin{itemize}
    \item Mensimulasikan gerak rotasi bumi pada porosnya.
    \item Memperagakan terjadinya siang dan malam.
    \item Menunjukkan letak garis lintang (ekuator/khatulistiwa) dan garis bujur secara akurat.
\end{itemize}

\subsection{Komponen Utama dan Simbolisasi pada Peta}
Agar peta dapat dibaca dan dimengerti oleh penggunanya, peta harus memiliki komponen-komponen baku.

\subsubsection{Arti Simbol-Simbol Peta}
Simbol peta adalah tanda atau gambar konvensional yang digunakan pada peta untuk mewakili keadaan atau objek sesungguhnya yang ada di lapangan. Simbol peta dikelompokkan menjadi empat jenis utama:
\begin{enumerate}
    \item \textbf{Simbol Titik:} Digunakan untuk menyajikan tempat atau data posisional, seperti ibu kota negara (simbol persegi dengan titik di tengah), gunung berapi (segitiga merah), gunung tidak aktif (segitiga hitam), dan pelabuhan (jangkar).
    \item \textbf{Simbol Garis:} Digunakan untuk menggambarkan objek yang memanjang. Contohnya adalah sungai (garis bercabang-cabang dan berkelok), jalan raya (garis lurus berwarna merah), rel kereta api (garis lurus dengan strip di tengahnya), serta batas provinsi atau batas negara.
    \item \textbf{Simbol Area/Luasan:} Digunakan untuk menggambarkan suatu area yang memiliki luasan tertentu, seperti danau, rawa, hutan, dan persawahan.
    \item \textbf{Simbol Warna:} Warna pada peta memberikan makna tertentu.
    \begin{itemize}
        \item \textit{Biru:} Mewakili wilayah perairan (laut, danau, sungai). Semakin gelap birunya, semakin dalam perairannya.
        \item \textit{Hijau:} Mewakili dataran rendah.
        \item \textit{Kuning:} Mewakili dataran tinggi atau perbukitan.
        \item \textit{Cokelat:} Mewakili wilayah pegunungan.
    \end{itemize}
\end{enumerate}

\subsubsection{Arti Skala dan Perhitungan Jarak}
Skala peta adalah angka perbandingan antara jarak pada peta dengan jarak sebenarnya di permukaan bumi. Skala menentukan tingkat kedetailan sebuah peta.
Terdapat tiga jenis skala: skala angka (misal 1:1.000.000), skala garis/grafis, dan skala verbal.

\textbf{Rumus Mencari Jarak Sebenarnya pada Peta:}
Secara matematis, hubungan antara Jarak Peta (JP), Skala (S), dan Jarak Sebenarnya (JS) dirumuskan sebagai berikut:
\[ \text{JS} = \text{JP} \times \text{Penyebut Skala} \]
Atau bisa juga ditulis:
\[ \text{JS} = \frac{\text{JP}}{\text{S}} \]

\textbf{Contoh Soal:}
Diketahui jarak antara kota A dan kota B pada peta adalah $5 \text{ cm}$. Peta tersebut memiliki skala $1 : 2.000.000$. Berapakah jarak sebenarnya antara kedua kota tersebut dalam satuan kilometer?
\begin{itemize}
    \item \textbf{Diketahui:} $\text{JP} = 5 \text{ cm}$, $\text{Penyebut Skala} = 2.000.000$
    \item \textbf{Penyelesaian:} 
    \[ \text{JS} = 5 \text{ cm} \times 2.000.000 = 10.000.000 \text{ cm} \]
    Karena $1 \text{ km} = 100.000 \text{ cm}$, maka:
    \[ \text{JS} = \frac{10.000.000}{100.000} \text{ km} = 100 \text{ km} \]
    \item \textbf{Jawaban:} Jarak sebenarnya adalah $100 \text{ km}$.
\end{itemize}

\subsection{Kondisi Geografis, Pembagian Waktu, dan Wilayah Indonesia}

\subsubsection{Pembagian Waktu di Indonesia (WIB, WITA, WIT)}
Letak astronomis Indonesia berada pada $6^\circ \text{ LU} - 11^\circ \text{ LS}$ dan $95^\circ \text{ BT} - 141^\circ \text{ BT}$. Karena panjang garis bujur Indonesia mencapai $46^\circ$, dan setiap $15^\circ$ garis bujur memiliki selisih waktu 1 jam, maka Indonesia dibagi menjadi tiga zona waktu berdasarkan Keputusan Presiden RI:
\begin{enumerate}
    \item \textbf{WIB (Waktu Indonesia Barat):} Berpatokan pada garis bujur $105^\circ \text{ BT}$. Memiliki selisih waktu $+7$ jam dari \textit{Greenwich Mean Time} (UTC+7).
    \item \textbf{WITA (Waktu Indonesia Tengah):} Berpatokan pada garis bujur $120^\circ \text{ BT}$. Memiliki selisih waktu $+8$ jam (UTC+8).
    \item \textbf{WIT (Waktu Indonesia Timur):} Berpatokan pada garis bujur $135^\circ \text{ BT}$. Memiliki selisih waktu $+9$ jam (UTC+9).
\end{enumerate}

\subsubsection{Jumlah Provinsi Terbaru dan Contoh Pembagian Wilayahnya}
Sejak era reformasi hingga tahun terbaru, terjadi banyak pemekaran wilayah di Indonesia. \textbf{Jumlah provinsi terbaru di Indonesia saat ini adalah 38 Provinsi}. Pemekaran terbaru difokuskan di wilayah Pulau Papua untuk mempercepat pemerataan pembangunan.

Berikut adalah contoh pembagian waktu berdasarkan provinsinya:
\begin{itemize}
    \item \textbf{Wilayah WIB:} Meliputi seluruh provinsi di Pulau Sumatera (Aceh, Sumatera Utara, Sumatera Barat, Riau, Kepri, Jambi, Bengkulu, Sumsel, Bangka Belitung, Lampung), seluruh provinsi di Pulau Jawa (Banten, DKI Jakarta, Jawa Barat, Jawa Tengah, DI Yogyakarta, Jawa Timur), serta dua provinsi di Kalimantan (Kalimantan Barat dan Kalimantan Tengah).
    \item \textbf{Wilayah WITA:} Meliputi sisa pulau Kalimantan (Kalimantan Selatan, Kalimantan Timur, Kalimantan Utara), seluruh provinsi di Pulau Sulawesi (Sulut, Gorontalo, Sulteng, Sulbar, Sulsel, Sultra), serta Bali, Nusa Tenggara Barat (NTB), dan Nusa Tenggara Timur (NTT).
    \item \textbf{Wilayah WIT:} Meliputi kepulauan Maluku (Maluku dan Maluku Utara) serta \textbf{seluruh enam provinsi di Pulau Papua} (Papua, Papua Barat, Papua Tengah, Papua Selatan, Papua Pegunungan, dan Papua Barat Daya).
\end{itemize}

\section{Interaksi Sosial, Nilai, Norma, dan Dinamika Sosial}

\subsection{Hakikat Interaksi Sosial}
Manusia sebagai \textit{Zoon Politicon} (makhluk sosial) tidak dapat hidup sendiri. Mereka membutuhkan orang lain untuk memenuhi kebutuhannya. Hubungan timbal balik, saling memengaruhi antara individu dengan individu, individu dengan kelompok, maupun kelompok dengan kelompok disebut sebagai \textbf{Interaksi Sosial}.

\textbf{Syarat Terjadinya Interaksi Sosial:}
Interaksi sosial tidak dapat terjadi begitu saja tanpa adanya dua syarat utama:
\begin{enumerate}
    \item \textbf{Kontak Sosial:} Pertemuan antara pihak-pihak yang berinteraksi. Kontak tidak harus bersentuhan fisik. Di era digital (seperti yang dipelajari siswa TKJ), kontak sosial bisa bersifat sekunder tidak langsung, seperti saling mengirim email, \textit{chatting} via WhatsApp, atau bermain \textit{game online} bersama.
    \item \textbf{Komunikasi:} Proses penyampaian pesan dari satu pihak kepada pihak lain yang diikuti dengan pemberian makna terhadap pesan tersebut.
\end{enumerate}

\subsection{Kelompok Sosial dan Pluralitas}
Interaksi manusia pada akhirnya akan membentuk \textbf{kelompok sosial}, yaitu kumpulan orang yang memiliki kesadaran bersama akan keanggotaan dan saling berinteraksi. 
Dalam masyarakat yang besar, timbul \textbf{pluralitas} (kemajemukan). Masyarakat majemuk terdiri atas berbagai macam suku, agama, ras, dan golongan profesi yang saling hidup berdampingan. Meskipun rawan gesekan, pluralitas merupakan kekayaan budaya yang dapat memperkuat ikatan nasional jika dikelola dengan sikap toleransi.

\subsection{Nilai, Norma, dan Sosialisasi}
Masyarakat agar bisa hidup tertib memerlukan pedoman. Pedoman ini diatur melalui konsep nilai dan norma.
\subsubsection{Pengertian Nilai dan Norma}
\begin{itemize}
    \item \textbf{Nilai Sosial:} Sesuatu yang dianggap baik, berharga, pantas, dan dicita-citakan oleh masyarakat. Contoh: Menghormati orang tua adalah nilai kebaikan.
    \item \textbf{Norma Sosial:} Wujud konkret dari nilai. Norma adalah aturan, patokan, atau ukuran yang mengikat tingkah laku manusia dalam masyarakat, serta dilengkapi dengan sanksi (hukuman) bagi pelanggarnya.
\end{itemize}

\subsubsection{Jenis-Jenis Norma Sosial}
\begin{enumerate}
    \item \textbf{Norma Agama:} Aturan yang bersumber dari wahyu Tuhan yang disampaikan melalui utusan-Nya. Sanksinya berupa dosa yang balasannya di akhirat, serta penderitaan batin.
    \item \textbf{Norma Kesusilaan:} Aturan yang bersumber dari hati nurani manusia mengenai apa yang benar dan apa yang salah. Contoh: Tidak berdusta. Sanksinya adalah rasa bersalah, penyesalan, atau malu.
    \item \textbf{Norma Kesopanan:} Aturan yang timbul dari kebiasaan pergaulan sekelompok masyarakat. Contoh: Menggunakan tangan kanan saat memberi sesuatu. Sanksinya berupa cemoohan, celaan, atau dikucilkan dari pergaulan.
    \item \textbf{Norma Hukum:} Aturan tertulis yang dibuat oleh lembaga negara yang berwenang (pemerintah). Bersifat mengikat dan memaksa. Sanksinya tegas dan nyata, seperti denda atau penjara.
\end{enumerate}

\subsubsection{Jenis Pelanggaran Norma Sekolah}
Sekolah merupakan miniatur masyarakat yang memiliki aturan tersendiri. Pelanggaran norma sekolah dapat menghambat proses kegiatan belajar mengajar. Contoh jenis pelanggarannya meliputi:
\begin{itemize}
    \item \textit{Pelanggaran Kedisiplinan:} Terlambat masuk kelas, membolos, pakaian seragam tidak sesuai atribut.
    \item \textit{Pelanggaran Etika/Kesopanan:} Berkata kasar kepada guru, melawan staf sekolah, tidak menyapa guru.
    \item \textit{Pelanggaran Hukum/Pidana (Kategori Berat):} Melakukan perundungan fisik maupun \textit{cyberbullying} (di media sosial), pencurian fasilitas sekolah (seperti kabel LAN, komponen PC di Lab TKJ), tawuran pelajar, atau mencontek saat ujian secara sistematis.
\end{itemize}

\subsubsection{Keteraturan Sosial}
Keteraturan sosial adalah kondisi di mana hubungan-hubungan sosial antaranggota masyarakat berlangsung selaras, serasi, dan stabil sesuai dengan nilai dan norma yang berlaku. Keteraturan ini tercipta melalui proses \textbf{Sosialisasi}, yaitu proses penanaman nilai dan norma dari satu generasi ke generasi berikutnya (melalui keluarga, sekolah, teman sebaya, dan media massa).

\subsection{Dinamika Sosial: Mobilitas, Perubahan, Konflik, dan Akomodasi}

\subsubsection{Mobilitas Sosial}
Mobilitas sosial adalah perpindahan posisi atau status sosial seseorang atau sekelompok orang dari satu lapisan ke lapisan lain dalam masyarakat.
\begin{itemize}
    \item \textbf{Mobilitas Vertikal Naik:} Seorang teknisi junior di perusahaan IT berhasil meraih sertifikasi Cisco dan diangkat menjadi Manajer Jaringan.
    \item \textbf{Mobilitas Vertikal Turun:} Seorang pengusaha ISP (Internet Service Provider) bangkrut dan menjadi karyawan biasa.
    \item \textbf{Mobilitas Horizontal:} Seorang guru SMK pindah mengajar dari sekolah A ke sekolah B dengan jabatan yang sama. Status sosialnya tidak berubah.
\end{itemize}

\subsubsection{Perubahan Sosial dan Budaya}
Perubahan sosial merupakan perubahan pada lembaga-lembaga kemasyarakatan yang memengaruhi sistem sosialnya, termasuk nilai, sikap, dan pola perilaku. Bagi siswa TKJ, memahami hal ini penting karena perkembangan teknologi merupakan motor utama perubahan sosial di era modern (revolusi industri 4.0). Perubahan gaya komunikasi masyarakat yang semula menggunakan surat menjadi pesan instan (\textit{chatting}) adalah salah satu bentuk nyata perubahan sosial budaya.

\subsubsection{Konflik, Kekerasan, dan Akomodasi}
Dinamika kehidupan sering memicu \textbf{konflik}, yaitu suatu proses sosial ketika individu atau kelompok berusaha memenuhi tujuannya dengan jalan menentang pihak lawan yang disertai dengan ancaman atau kekerasan. Jika konflik gagal diselesaikan, ia akan memuncak menjadi \textbf{kekerasan} fisik maupun mental.

Untuk meredakan konflik agar masyarakat tidak terpecah belah, dilakukan upaya \textbf{Akomodasi}. Akomodasi adalah proses penyesuaian diri individu atau kelompok untuk meredakan ketegangan tanpa menghancurkan pihak lawan.

\textbf{Cara Menanggulangi Konflik (Bentuk-bentuk Akomodasi):}
\begin{enumerate}
    \item \textbf{Kompromi:} Masing-masing pihak yang berkonflik saling mengurangi tuntutannya demi mencapai kesepakatan.
    \item \textbf{Mediasi:} Melibatkan pihak ketiga sebagai penengah, namun pihak ketiga tersebut bersikap netral dan \textit{hanya sebagai penasihat} (tidak memiliki hak memutuskan).
    \item \textbf{Arbitrasi:} Melibatkan pihak ketiga yang memiliki otoritas lebih tinggi. Keputusan pihak ketiga bersifat \textit{mengikat dan memaksa} kedua belah pihak.
    \item \textbf{Konsiliasi:} Usaha mempertemukan pimpinan dari pihak-pihak yang berkonflik dalam suatu perundingan resmi.
    \item \textbf{Ajudikasi:} Penyelesaian konflik melalui jalur hukum (pengadilan).
\end{enumerate}

\section{Ilmu Ekonomi Dasar: Kebutuhan Manusia, Pelaku Ekonomi, dan Kelangkaan}

Ilmu ekonomi mempelajari bagaimana manusia memilih cara menggunakan sumber daya alam yang terbatas untuk memenuhi kebutuhannya yang tidak terbatas.

\subsection{Karakteristik Pemenuhan Kebutuhan Manusia}
\textbf{Karakter pemenuhan kebutuhan manusia memiliki ciri-ciri dasar:}
\begin{itemize}
    \item \textbf{Tidak Terbatas:} Manusia pada dasarnya tidak pernah merasa puas. Jika satu kebutuhan terpenuhi, akan muncul tuntutan kebutuhan lain yang lebih tinggi.
    \item \textbf{Dinamis (Selalu Berubah):} Kebutuhan terus berkembang seiring berjalannya waktu, bertambahnya usia, tingkat pendapatan, lingkungan geografis, dan kemajuan peradaban/teknologi. 
    \item Di sisi lain, sumber daya untuk memenuhinya (alat pemuas kebutuhan) bersifat \textbf{Terbatas}. Jurang pemisah antara kebutuhan yang tak terbatas dengan sumber daya yang terbatas ini menghasilkan masalah ekonomi utama, yaitu \textbf{Kelangkaan (Scarcity)}.
\end{itemize}

\subsection{Pengklasifikasian Kebutuhan Manusia}
Kebutuhan manusia dapat dibagi berdasarkan beberapa kategori, yang paling sering keluar dalam ujian adalah berdasarkan intensitas dan sifatnya.

\subsubsection{Kebutuhan Berdasarkan Intensitas (Tingkat Kepentingan)}
\begin{enumerate}
    \item \textbf{Kebutuhan Primer (Pokok):} Kebutuhan mutlak yang harus dipenuhi pertama kali agar manusia bisa bertahan hidup. Jika tidak terpenuhi, kelangsungan hidup manusia terancam. Contoh: Makanan bergizi, pakaian (sandang), dan tempat tinggal (papan/rumah).
    \item \textbf{Kebutuhan Sekunder (Pelengkap):} Kebutuhan yang dipenuhi setelah kebutuhan primer tercukupi. Kebutuhan ini bertujuan untuk menambah kenyamanan hidup. Contoh: Televisi, kasur empuk, sepeda motor standar, kulkas, dan \textit{smartphone} kelas menengah.
    \item \textbf{Kebutuhan Tersier (Mewah):} Kebutuhan akan barang-barang mewah yang dipenuhi untuk menaikkan prestise atau status sosial di mata masyarakat. Pemenuhan kebutuhan ini biasanya dilakukan oleh kelompok ekonomi menengah ke atas. Contoh: Mobil \textit{sport}, perhiasan berlian, liburan ke luar negeri, tas bermerek (\textit{branded}).
\end{enumerate}

\subsubsection{Kebutuhan Berdasarkan Sifatnya (Materil dan Imateril)}
\begin{enumerate}
    \item \textbf{Kebutuhan Materil (Jasmani):} Kebutuhan yang berhubungan dengan wujud fisik badan manusia. Barang yang dibutuhkan berwujud nyata, dapat dilihat, dan diraba. Kebutuhan ini bertumpu pada pelestarian fisik. Contoh: Nasi untuk makan, obat-obatan saat sakit, pakaian agar tidak kedinginan, alat peraga praktik TKJ (seperti kabel UTP, \textit{Router Mikrotik}, komputer).
    \item \textbf{Kebutuhan Imateril (Rohani):} Kebutuhan yang bersifat tidak berwujud fisik, melainkan menyangkut ketenangan jiwa, batin, dan emosional manusia. Pemenuhan kebutuhan ini memberikan kepuasan secara psikologis. Contoh:
    \begin{itemize}
        \item Menjalankan ibadah sesuai agama masing-masing.
        \item Menonton konser musik atau bermain \textit{game} untuk hiburan/rekreasi.
        \item Menuntut ilmu di sekolah dan membaca buku (pendidikan).
        \item Mendapatkan perhatian, kasih sayang keluarga, serta rasa aman dari gangguan kejahatan.
    \end{itemize}
\end{enumerate}

\subsection{Pelaku Ekonomi dan Peranannya}
Kegiatan perekonomian (produksi, distribusi, konsumsi) tidak berjalan sendiri, melainkan digerakkan oleh subjek-subjek ekonomi. Terdapat empat kelompok pelaku ekonomi utama dalam suatu negara, yaitu:

\begin{enumerate}
    \item \textbf{Rumah Tangga Konsumsi (RTK) / Konsumen:} 
    Individu atau kelompok masyarakat yang mengonsumsi barang dan jasa untuk memenuhi kebutuhan hidupnya. RTK memiliki peran ganda: selain sebagai pembeli produk, mereka juga merupakan pemilik faktor-faktor produksi (seperti tenaga kerja, tanah, dan modal) yang ditawarkan kepada perusahaan. Atas jasanya tersebut, RTK menerima imbalan berupa upah, uang sewa, atau dividen/bunga.

    \item \textbf{Rumah Tangga Produksi (RTP) / Perusahaan:} 
    Organisasi ekonomi yang didirikan dengan tujuan menghasilkan barang dan jasa. Perusahaan menyewa/membeli faktor produksi dari RTK, mengolahnya, dan menjual hasilnya kembali ke masyarakat luas guna mendapatkan laba (keuntungan). RTP dituntut untuk menerapkan prinsip efisiensi dan inovasi agar produknya laku di pasar.

    \item \textbf{Rumah Tangga Negara (RTN) / Pemerintah:} 
    Pelaku ekonomi yang memiliki kekuasaan secara hukum untuk mengatur dan menyehatkan jalannya perekonomian negara. Peran pemerintah antara lain:
    \begin{itemize}
        \item \textit{Sebagai Pengatur/Regulator:} Membuat undang-undang, menetapkan tarif pajak, memberikan izin usaha, dan menjaga kestabilan harga.
        \item \textit{Sebagai Produsen:} Mengelola sektor strategis yang menguasai hajat hidup orang banyak melalui BUMN (Contoh: PLN mengelola listrik, Pertamina mengelola minyak tanah, PT KAI mengelola kereta api).
        \item \textit{Sebagai Konsumen:} Pemerintah juga butuh berbelanja untuk operasional negara, seperti membeli kendaraan dinas, alat tulis kantor untuk dinas, dan senjata untuk TNI/Polri.
    \end{itemize}

    \item \textbf{Masyarakat Luar Negeri:} 
    Di era modern, tidak ada negara yang bisa berdiri sendiri (Autarki). Negara selalu membutuhkan interaksi dengan dunia luar. Peran masyarakat luar negeri meliputi aktivitas perdagangan internasional (\textbf{Ekspor} barang keluar untuk mendatangkan devisa, dan \textbf{Impor} barang dari luar untuk memenuhi kebutuhan yang tidak bisa diproduksi di dalam negeri), pertukaran tenaga kerja, serta penanaman modal (investasi asing).
\end{enumerate}

\end{document}
